% ============================================================================
%  Problems: จำนวนจริง (Real Numbers)
%  Ordered from easy (basic) → hard (POSN level).
%  Shared by exercise.tex and answer-key.tex
% ============================================================================

% --- Part 1: พื้นฐาน (Basic) ---
\examsection{ตอนที่ 1: พื้นฐาน — เศษส่วน ทศนิยม และระบบจำนวน}

\problem{
  จงระบุว่าจำนวนต่อไปนี้เป็นจำนวนตรรกยะหรืออตรรกยะ
  \begin{enumerate}
    \item[(ก)] $\frac{3}{7}$
    \item[(ข)] $\sqrt{5}$
    \item[(ค)] $0.75$
    \item[(ง)] $\pi$
  \end{enumerate}
}{
  (ก) ตรรกยะ \quad (ข) อตรรกยะ \quad (ค) ตรรกยะ \quad (ง) อตรรกยะ
}

\problem{
  จงเขียน $0.333\ldots$ ในรูปเศษส่วนอย่างต่ำ
}{
  ให้ $x = 0.333\ldots$,
  $10x = 3.333\ldots$,
  $9x = 3$,
  $x = \frac{1}{3}$
}

\problem{
  จงคำนวณค่าต่อไปนี้และตอบในรูปเศษส่วนอย่างต่ำ
  \begin{enumerate}
    \item[(ก)] $\dfrac{2}{5} + \dfrac{1}{3}$
    \item[(ข)] $\dfrac{3}{4} \times \dfrac{2}{5}$
    \item[(ค)] $\dfrac{7}{8} \div \dfrac{3}{4}$
  \end{enumerate}
}{
  (ก) $\frac{11}{15}$

  (ข) $\frac{3}{10}$

  (ค) $\frac{7}{6} = 1\frac{1}{6}$
}

\problem{
  จงเขียนจำนวนต่อไปนี้ในรูปร้อยละ
  \begin{enumerate}
    \item[(ก)] $0.35$
    \item[(ข)] $\frac{3}{5}$
    \item[(ค)] $1.2$
  \end{enumerate}
}{
  (ก) $35\%$ \quad (ข) $60\%$ \quad (ค) $120\%$
}

\problem{
  จงเขียนอัตราส่วนต่อไปนี้ในรูปอย่างต่ำ
  \begin{enumerate}
    \item[(ก)] $12 : 18$
    \item[(ข)] $0.5 : 2$
  \end{enumerate}
}{
  (ก) $2 : 3$ \quad (ข) $1 : 4$
}

% --- Part 2: การประยุกต์ (Intermediate) ---
\examsection{ตอนที่ 2: การประยุกต์ — สมบัติ สัดส่วน และค่าสัมบูรณ์}

\problem{
  จงระบุว่าแต่ละข้อต่อไปนี้ใช้สมบัติใดของระบบจำนวนจริง
  \begin{enumerate}
    \item[(ก)] $3 + 5 = 5 + 3$
    \item[(ข)] $(2 \times 3) \times 4 = 2 \times (3 \times 4)$
    \item[(ค)] $2 \times (3 + 5) = 2 \times 3 + 2 \times 5$
    \item[(ง)] $7 + 0 = 7$
    \item[(จ)] $5 \times \frac{1}{5} = 1$
  \end{enumerate}
}{
  (ก) Commutative (การสลับที่การบวก)

  (ข) Associative (การเปลี่ยนหมู่การคูณ)

  (ค) Distributive (การแจกแจง)

  (ง) Identity (เอกลักษณ์การบวก)

  (จ) Inverse (อินเวอร์สการคูณ)
}

\problem{
  ถ้าขนมปัง 5 ชิ้น ราคา 75 บาท
  ขนมปัง 8 ชิ้นจะราคากี่บาท?
  (ใช้สัดส่วนตรง)
}{
  $\frac{5}{75} = \frac{8}{x}$,
  $5x = 600$,
  $x = 120$ บาท
}

\problem{
  คนงาน 6 คน ทำงานหนึ่งเสร็จใน 10 วัน
  ถ้าใช้คนงาน 15 คน จะทำงานเสร็จในกี่วัน?
  (ถือว่าทุกคนทำงานได้เท่ากัน ใช้สัดส่วนผกผัน)
}{
  $6 \times 10 = 15 \times x$,
  $x = 4$ วัน
}

\problem{
  จงแก้สมการต่อไปนี้
  \begin{enumerate}
    \item[(ก)] $|x| = 8$
    \item[(ข)] $|x - 4| = 3$
  \end{enumerate}
}{
  (ก) $x = 8$ หรือ $x = -8$

  (ข) $x = 7$ หรือ $x = 1$
}

\problem{
  จงแก้อสมการ $|x + 2| < 5$ และเขียนคำตอบในรูปช่วง
}{
  $-5 < x + 2 < 5$,
  $-7 < x < 3$,
  $x \in (-7, 3)$
}

\problem{
  จงหาค่าต่อไปนี้ในรูปอย่างง่าย
  \begin{enumerate}
    \item[(ก)] $2^3 \times 2^4$
    \item[(ข)] $\dfrac{3^5}{3^2}$
    \item[(ค)] $\sqrt{50}$ (เขียนในรูป $a\sqrt{b}$)
  \end{enumerate}
}{
  (ก) $2^7 = 128$

  (ข) $3^3 = 27$

  (ค) $\sqrt{50} = \sqrt{25 \times 2} = 5\sqrt{2}$
}

% --- Part 3: ระดับ สอวน. (POSN Level) ---
\examsection{ตอนที่ 3: ระดับข้อสอบ สอวน.}

\problem{
  กำหนดตัวดำเนินการใหม่ $a \diamond b = a + b - ab$
  จงตรวจสอบว่า $\diamond$ มีสมบัติต่อไปนี้หรือไม่
  \begin{enumerate}
    \item[(ก)] สมบัติการสลับที่ (Commutative)
    \item[(ข)] สมบัติการเปลี่ยนหมู่ (Associative)
  \end{enumerate}
  พร้อมทั้งหาค่าเอกลักษณ์ของ $\diamond$ (ถ้ามี)
}{
  (ก) $a \diamond b = a + b - ab$, $b \diamond a = b + a - ba$ — เหมือนกัน

  \textbf{มี}สมบัติการสลับที่

  (ข) $(a \diamond b) \diamond c = (a+b-ab) + c - (a+b-ab)c$
  $= a+b+c-ab-ac-bc+abc$

  $a \diamond (b \diamond c) = a + (b+c-bc) - a(b+c-bc)$
  $= a+b+c-bc-ab-ac+abc$

  เท่ากัน — \textbf{มี}สมบัติการเปลี่ยนหมู่

  เอกลักษณ์ $e$: $a \diamond e = a \Rightarrow a + e - ae = a \Rightarrow e(1-a) = 0$

  สำหรับทุก $a$: $e = 0$

  ดังนั้นเอกลักษณ์คือ \textbf{0}
}

\problem{
  จงเขียน $0.\overline{27}$ (ทศนิยมซ้ำ 27) ในรูปเศษส่วนอย่างต่ำ
}{
  ให้ $x = 0.272727\ldots$

  $100x = 27.272727\ldots$

  ลบกัน: $100x - x = 27$,
  $99x = 27$,
  $x = \frac{27}{99} = \frac{3}{11}$
}

\problem{
  จงหาค่า $x$ ทั้งหมดที่สอดคล้องกับสมการ $|2x - 1| = |x + 3|$
}{
  ยกกำลังสองทั้งสองข้าง: $(2x-1)^2 = (x+3)^2$

  $4x^2 - 4x + 1 = x^2 + 6x + 9$

  $3x^2 - 10x - 8 = 0$

  $(3x + 2)(x - 4) = 0$

  $x = -\frac{2}{3}$ หรือ $x = 4$

  \textbf{ตรวจสอบ:} ทั้งสองค่าใช้ได้ $\checkmark$
}

\problem{
  กำหนด $x$ และ $y$ เป็นจำนวนจริงบวก โดยที่ $\sqrt{x\sqrt{y}} = 4$
  และ $\sqrt{y\sqrt{x}} = 2$ จงหาค่าของ $x$ และ $y$
}{
  $\sqrt{x\sqrt{y}} = 4 \Rightarrow x\sqrt{y} = 16 \Rightarrow \sqrt{y} = \frac{16}{x}$

  $\sqrt{y\sqrt{x}} = 2 \Rightarrow y\sqrt{x} = 4$

  แทน $\sqrt{y}$: $(\frac{16}{x})^2 \cdot \sqrt{x} = 4$

  $\frac{256}{x^2} \cdot \sqrt{x} = 4$
  $\Rightarrow 64 = x^2 \cdot x^{-1/2}$
  $\Rightarrow 64 = x^{3/2}$
  $\Rightarrow x = 16$

  จาก $\sqrt{y} = \frac{16}{x} = 1$
  $\Rightarrow y = 1$

  \textbf{คำตอบ:} $x = 16$, $y = 1$
}

\problem{
  (ข้อสอบจริงปี 65 ข้อ 6) จงหาค่าของ $k$ และ $m$ ที่สอดคล้องกับสมการ $\dfrac{x^{2k} \times x^m}{x^m \div x^k}$ ตามลำดับ
}{
  $k=15, m$ มีค่าเป็นเท่าไหร่ก็ได้ที่ไม่ทำให้ $x^m=0$
}

\problem{
  (ข้อสอบจริงปี 65 ข้อ 34) ให้ $a$ เป็นจำนวนจริงใด ๆ และ $b$ เป็นผกผันการคูณของ $a$ 
  ถ้า $c$ เป็นจำนวนจริง ซึ่ง $2a^2+abc=3$ แล้วข้อใดถูกต้อง 
  \begin{enumerate}
    \item[ก.] $c$ มีค่าเท่ากับเอกลักษณ์การบวก
    \item[ข.] $c$ มีค่าเท่ากับเอกลักษณ์การคูณ
    \item[ค.] $c$ มีค่าเท่ากับผกผันการคูณของ $a$
    \item[ง.] $c$ มีค่าเท่ากับผกผันการคูณของ $b$
  \end{enumerate}
}{
  ข.
}

\problem{
  (ข้อสอบจริงปี 68 ข้อ 26) มีพนักงานสามคนช่วยกันแพ็คสินค้า ได้แก่ A, B, C สำหรับสินค้าจำนวนหนึ่ง
  \begin{itemize}
    \item ถ้า A และ B ช่วยกันจนเสร็จ ใข้เวลา 18 วัน
    \item ถ้า B และ C ช่วยกันจนเสร็จ ใข้เวลา 27 วัน
    \item ถ้า C และ A ช่วยกันจนเสร็จ ใข้เวลา 36 วัน
  \end{itemize}
  ถ้าสามคนช่วยกันแพ็คมา 8 วัน แล้ว A และ C ลา B ต้องแพ็คที่เหลือทั้งหมด ต้องใช้เวลากี่วันในการแพ็คส่วนที่เหลือให้เสร็จ
}{
  $14$ วัน
}
